% Generated draft: requires mathematical and source review.
\documentclass[11pt]{article}
\usepackage[margin=25mm]{geometry}
\usepackage{fontspec}
\setmainfont{Latin Modern Roman}
\usepackage{amsmath,amssymb,mathtools,unicode-math}
\setmathfont{Latin Modern Math}
\usepackage{enumitem,needspace,xurl,hyperref}
\setlength{\parindent}{0pt}
\setlength{\parskip}{5pt}
\setlength{\emergencystretch}{3em}
\newcommand{\N}{\mathbb{N}}
\newcommand{\Z}{\mathbb{Z}}
\newcommand{\Q}{\mathbb{Q}}
\newcommand{\R}{\mathbb{R}}
\newcommand{\C}{\mathbb{C}}
\newcommand{\F}{\mathbb{F}}
\newcommand{\A}{\mathbb{A}}
\newcommand{\PP}{\mathbb{P}}
\newcommand{\E}{\mathbb{E}}
\newcommand{\eps}{\varepsilon}
\newcommand{\dd}{\,\mathrm{d}}
\begin{document}
\section*{UnsolvedMath ID: 30001261}
\textbf{Source problem code:} catalogue entry
\subsection{Definitions and mathematical statement}
A group \(G\) is virtually cyclic if it has a cyclic subgroup \(C\leq G\) of finite index, meaning \([G:C]<\infty\). The cyclic subgroup is allowed to be trivial.

Under the standard interpretation of \(G\) as a discrete group, a model for the free classifying space \(EG\) is a contractible free \(G\)-CW complex \(X\). Freeness means that every point stabilizer is trivial. Such a complex is finite as a \(G\)-CW complex if it has finitely many \(G\)-orbits of cells. This differs from requiring its underlying ordinary CW complex to have finitely many cells.

The supplied question is: if a group \(G\) admits a finite model for its free classifying space \(EG\), must \(G\) be virtually cyclic? The record does not specify which meaning of finite is intended. Under the two conventions above, the respective candidate statements are
\[
\begin{aligned}
&\text{For every discrete group }G,\\
&\quad \bigl(\exists\text{ a contractible free }G\text{-CW complex with finitely many }G\text{-orbits of cells}\bigr)\\
&\qquad\Longrightarrow \bigl(\exists C\leq G:\ C\text{ is cyclic and }[G:C]<\infty\bigr),
\end{aligned}
\]
and
\[
\begin{aligned}
&\text{For every discrete group }G,\\
&\quad \bigl(\exists\text{ a contractible free }G\text{-CW complex whose underlying CW complex is finite}\bigr)\\
&\qquad\Longrightarrow \bigl(\exists C\leq G:\ C\text{ is cyclic and }[G:C]<\infty\bigr).
\end{aligned}
\]
These are alternative readings, not assertions that both express the intended problem. The accompanying attribution to the Juan--Pineda--Leary conjecture leaves an additional unresolved issue concerning the intended classifying space. The supplied evidence does not justify replacing the explicitly stated free space \(EG\) by a classifying space for a family of subgroups. A unique intended formulation therefore requires source review.
\subsection{Short English statement}
The record asks whether a group with a finite free classifying space must contain a cyclic subgroup of finite index. Both the meaning of finite and the compatibility of the stated classifying space with the accompanying conjecture attribution require clarification. The displayed formulations preserve these unresolved choices.
\subsection{Sources}
\begin{itemize}
\item Supplied local registry record: \texttt{data/unsolvedmath/problems.json}, permanent UnsolvedMath ID \(30001261\), source problem code \texttt{catalogue entry}, title \emph{Virtual Cyclicity from Finite Free Classifying Spaces}. The supplied collection is \emph{Oberwolfach Reports Open Problems}; the proposed year is 2009, and no individual proposer is supplied.
\item Supplied report citation: \emph{Manifold Perspectives} (2009), \emph{Oberwolfach Reports}, pp.\ 1487--1546, DOI 10.4171/owr/2009/27. This is also the literature-source entry labeled ``original OWR report'': \url{https://doi.org/10.4171/owr/2009/27}.
\item Supplied literature-source entry labeled ``Juan--Pineda--Leary'': \url{https://eprints.soton.ac.uk/199397/}.
\item Supplied literature-source entry labeled ``current status survey'': \url{https://eprints.soton.ac.uk/420131/2/hierarchical_F_final.pdf}.
\item Literature assessment dated exactly \texttt{2026-08-21}: the supplied curation labels the Juan--Pineda--Leary conjecture open in general and reports that it is proved for several classes, including elementary amenable and hyperbolic groups. This is supplied dated curation derived from the OWR material, not an independently verified assessment. Its applicability to the free-classifying-space wording above remains unresolved; time-sensitive status and cited claims require independent verification.
\end{itemize}
\end{document}
